\subsection{The actual return of the reduced type III family}

\begin{statement}{P3K.24}{Proved}
For the family of P3F.37, put $m=4l+3$, $l\ge0$, and retain
its actual parameter $\lambda=r^m$. Its positive $r$-return
satisfies
\[
 A(\Phi)=-(3l+2),\qquad
 \mu(\Sigma_\Phi)=-\frac{3l+2}{28}\pmod{\mathbb Z},
 \qquad [\Phi]=26-3l\pmod{28}.                         \tag{K24.1}
\]
The same construction and formula hold for the explicit source
$y^2=x^3+tx-t$, $G=(1,1)$, with fibres III*$+$II$+I_1$,
using its $I_1$ as the sole simple whole-fibre zero. As $l$ runs
from zero to twenty-seven, either set of families realizes all
oriented classes in $\Theta_7$. For a nonzero residue its mapping
torus is not diffeomorphic to $S^1\times S^6$.
\end{statement}
\begin{proof}
\par\medskip\noindent\textbf{Physical characters and the closing arrow.}
First use the source of P3F.37. Write $t_c$ for the selected cusp,
$q=t-t_c$, $s=\log q/(2\pi i)$, and $v=\log r/(2\pi i)$.
Use the physical source Abel coordinate $a=\zeta_1$ and the
physical line character $V=e^{2\pi i\zeta_2}$. Let $g(q)$
be the holomorphic Tate logarithm of $+G$, and let $C(q)$ be
the bounded unit logarithm in the linearization. The actual maps
are
\[
 \begin{aligned}
 B(a,V)&=(a-g,e^{2\pi iC}qr^mV),\\
 B^{-1}(a,V)&=(a+g,e^{-2\pi iC}q^{-1}r^{-m}V).
 \end{aligned}
\]
Thus the positive outward suspension loop is the period
\[
 \gamma_{\rm out}=\binom{g}{-C-s-mv},\qquad
 \delta=\binom01.                                    \tag{K24.2}
\]
Delta is positive physical scalar rotation. In particular the
period (K24.2) is $B^{-1}$, the same outward loop as in the
cap meridian and the Seifert marking; neither coordinate has
been silently inverted.

The annular arrow in P3F.37 is
$F_c:\text{ordinary}\to\text{cap}$, where
$c=\log\lambda/(2\pi i)$. Its Lie translation satisfies
$g_c(s_\infty+1)=Dg_c(s_\infty)+a_c$, so its scalar variation
is $(c-c_0)s_\infty\delta/m$. The cap is identical at $c+m$.
Consequently the inverse closing is $F_c^{-1}F_{c+m}$, with
multiplier $q_\infty$ at infinity, and globally
$(t-t_c)^{-1}$. The positive parameter deck is the inverse:
\[
              (v,\zeta)\longmapsto(v+1,\zeta+s\delta). \tag{K24.3}
\]
This square fixes its direction, including its bounded unit;
it does not infer the direction from a finite-difference identity.

On the auxiliary cusp cover $q=Q_0^{2m}$ take the extending origins
\[
 o(s,v)=\binom{-gs/m}{Cs/m+s^2/(2m)+sv},\qquad
 \sigma(s,v)=\binom{gv}{-Cv-m(v^2-v)/2}.
\]
Their exact differences are
\[
 \begin{aligned}
 o(s,v+1)-o(s,v)&=s\delta,\\
 \sigma(s,v+1)-\sigma(s,v)&=s\delta+\gamma_{\rm out},\\
 o(s+2m,v)-o(s,v)&=-2\gamma_{\rm out}+2m\delta,\\
 \sigma(s+2m,v)-\sigma(s,v)&=0.
 \end{aligned}                                                    \tag{K24.4}
\]
The factor two in this auxiliary cover is necessary for these
particular origins: $m$ is odd, and the one-$m$ shift would
contain $m\delta/2$. Every discrepancy in (K24.4) is an actual
integral period. The cusp origin has bounded characters in $q$
and extends in the unimodular cusp charts. Subtracting $o$
removes the parameter deck on the ordinary collar. For fixed
nonzero $q$, its remaining multiplier is a unit times $r^m$;
the actual $m$-node chain
$X_i=V/r^i$, $Y_i=r^{i+1}/V$, $X_iY_i=r$ gives its extending
origin. These are the two specified collar comparisons, without
an asserted two-variable proper analytic corner.

Their difference and angular change are
\[
 F=\sigma-o=
 \binom{g(s/m+v)}{-C(s/m+v)-(s+mv)^2/(2m)+mv/2},
\]
\[
 F(s,v+\theta)-F(s,v)
 =\theta\gamma_{\rm out}
          -\frac m2\theta(\theta-1)\delta.             \tag{K24.5}
\]
The last term is an endpoint-fixed periodic unitary gauge.
It vanishes in the circle quotient. The whole displayed angular
difference changes under $s\mapsto s+2m$ by $-2m\theta\delta$, so its retained
line transition has winding $-2m$. The actual positive $r$-loop
therefore makes one positive outward gamma turn. The auxiliary
cusp cover has not changed the parameter-loop count.

\par\medskip\noindent\textbf{The actual Seifert space and page.}
The cap marking is
\[
 (z,u)\sim(\zeta_mz,u-1),\qquad q_\infty=z^m,
 \qquad u=m\gamma_{\rm out}-\arg(q_\infty)/(2\pi).
\]
Let $R(z)>0$ solve $R^2+|z|^2R^{2/m}=1$. The explicit map
\[
 (z,u)\longmapsto
 (z_1,z_2)=
 \bigl(zR^{1/m}e^{2\pi iu/m},R e^{2\pi iu}\bigr)       \tag{K24.6}
\]
identifies the marked three-dimensional quotient with the
weighted circle fibration on $S^3$ of weights $(1,m)$.
The root and phase factors cancel under its indicated equivalence.
Conversely a real argument of $z_2$ recovers $u$ and then $z$;
another argument gives exactly that equivalence. Its base
coordinate is $t=z_2/z_1^m$, and its positive gamma rotation
is $(z_1,z_2)\mapsto(e^{2\pi i\theta}z_1,
e^{2\pi im\theta}z_2)$. In particular it retains
$q\mapsto-h$, $\gamma\mapsto mh$.

Choose a regular binding $t=t_*$ near, but distinct from,
$t_c$. Its positive Milnor function is $z_2-t_*z_1^m$.
The affine page $b-t_*a^m=1$ embeds by
$(a,b)\mapsto(\rho a,\rho^mb)$, where the unique positive
$\rho$ solves $\rho^2|a|^2+\rho^{2m}|b|^2=1$.
The inverse uses the positive weighted root of the Milnor
function. Its projection is exactly
\[
                         t=t_*+a^{-m}.                 \tag{K24.7}
\]
The compactification is $\mathbb P^1_a$, branched to order
$m$ at $a=0$ over the additive cap and at $a=\infty$ over
the regular binding. Its open page $C^\circ$ is a disc with
one boundary circle. The chosen cusp and the other finite
critical values each have $m$ unramified preimages. This is
the regular-binding cover; no cusp-branched map is substituted.

\par\medskip\noindent\textbf{The root slice, torsor and line.}
The full cap orbit is the actual mapping torus of $a_0$ on
the resolved III source. In (K24.6) the page equation is
\[
 2\pi u+\arg(1-t_*z^m)=\varphi\pmod{2\pi}.
\]
On a small cap disc choose its real argument. The lifted
graphs are $u=u_0(z)+j$, $j\in\mathbb Z$, where
$u_0=\varphi/(2\pi)-\arg(1-t_*z^m)/(2\pi)$.
The actual $A_r$ moves adjacent graphs and applies $a_0$ to
the source point. Every orbit has a unique representative
on $j=0$: $u_0(\zeta_mz)=u_0(z)$, so an $A_r^j$ equality
on that graph forces $j=0$. Thus
\[
                         p\longmapsto[p,u_0(z(p))]
\]
identifies the entire resolved III source with the actual page
near that core, including both components and their tangency
point. No residual stabilizer quotient remains. The delta
bundle on the graph is the norm circle of the actual line
$N=\mathcal O(2G'-O)$ with a smooth positive radius. Radial
normalization identifies it with the unit circle of $N$.

The other finite page pieces are unramified source pullbacks.
At the cusp use P3F.9 relative to its gamma and source boundary
marking. The degree-$m$ page can avoid the short gamma arc of
the fixed-sphere compression exactly as in P3K.6. This gives
an actual affine elliptic fibration over $C^\circ$, locally
translation-isomorphic to the minimal pullback under (K24.7).

The comparison can fix the binding collar. In one sector of
the cover the three original nodal vanishing paths lift with
their source marking. Their primitive vanishing cycles span
the full elliptic lattice, by P3F.37's trace-modulo-$d^2$
argument. Thus the global lattice coinvariant group is zero.
For the period local system $\Lambda$ on the punctured page,
compact-support duality gives $H_c^2(\Lambda)=H_0(\Lambda)=0$.
The quotient $j_*\Lambda/j_!\Lambda$ is a finite skyscraper
sheaf with no $H^1$, so
\[
             H^2(C^\circ,\partial C^\circ;j_*\Lambda)=0. \tag{K24.8}
\]
The smooth extending Neron translation sheaf has an identity
part with exponential kernel $j_*\Lambda$ and a fine
Lie-algebra sheaf. At III the identity group is additive;
the local exponential surjectivity follows from its derivative.
At a node it is the multiplicative exponential. Its component
quotient is the skyscraper $\mathbb Z/2$ at III, with zero
$H^1$ relative to the ordinary binding collar. Remove the
component cocycle by local component sections, lift the remainder
to the Lie bundle, remove its integral two-cocycle using (K24.8),
and solve the remaining Lie cocycle by a partition of unity
vanishing on that collar. The resulting translations give
mutually inverse smooth maps to the sectioned pullback, equal
to the prescribed comparison at the binding. Their extensions
over III are the actual extending source translations.

For the line comparison use the rigidified Picard map before
pulling back by these smooth translations. On a holomorphic
translation chart form $\operatorname{Hom}(N,T^*N)$ using
the universal extending source translation $T$. Its restriction
to each compact fibre is trivial. On regular elliptic fibres
this is the degree-zero group law; on the III fibre $N$ is
trivial, as the doubling calculation below shows; on a node
the normalization's multiplicative translation preserves its
gluing unit. The retained II case below has normalization
coordinate $u$ translated by a constant: its cusp coordinate
$w=1/u$ changes to $w/(1+tw)$, preserving the first jet
classifying its degree-zero Picard line. Thus the fibre $H^0$
is always one-dimensional. Cohomology and base change makes
the direct image a line, and its evaluation is an isomorphism
on each fibre. Evaluation along $O$ rigidifies it and identifies
it with a base line. Pulling this actual isomorphism back along
the smooth translation gives a smooth, fibrewise holomorphic
line comparison; it does not assign a holomorphic divisor to
a merely smooth total-space map.

The actual page line consequently has Chern class that of
$P-O$, $P=2G'$, up to a base pullback. Its absolute base class
vanishes on $C^\circ$. Retaining the prescribed boundary
gauge can leave an integral relative class $nF$, whose square
and product with $M$ vanish because $F^2=FM=0$. On the marked
boundary $T^3$ the degree-zero line and its commuting
translation lift admit a flat representative, so the boundary
Chern class is zero. The relative cup-product calculation
P3K.5 preserves the square for all such retained gauges.

\par\medskip\noindent\textbf{Characteristic numbers and the signed Thom map.}
The compact pullback has one III and $3m$ nodal fibres.
Hence
\[
 \chi=\frac{m+1}{4}=l+1,\qquad \sigma=-8\chi.
\]
The pulled section $P=2G'$ has height $2m$, identity component
at III, and no other reducible-fibre contribution. Thus
\[
 P\cdot O=m-\chi=3l+2,\qquad
 (P-O)^2=-2\chi-2P\cdot O=-2m.                         \tag{K24.9}
\]
These intersections have an explicit location. For (F37.1)
the doubling slope at $G'$ is $1/(2z^{3l+2})$, and the
zero-section parameter of $2G'$ is $2z^{3l+2}$ times a unit.
All $3l+2$ intersections occur at this actual cap; none is
inserted at the regular binding. In particular the central
degree-zero line there is trivial, as used above.

The regular elliptic cap and the real compression cap of
P3K.6 have zero signature. The latter kills the nodal vanishing
cycle, and its $H^2$ restriction to the marked $T^3$ is
injective. Thus the actual page gives
\[
                 s_\delta=-2m,\qquad p_{\rm rel}=-24\chi.
\]
We now retain the physical orientation reversal as well as
the Thom meridian. With $V=\rho e^{i\phi}$ and positive
delta $\partial_\phi$, the complex orientation is
$\operatorname{or}(\text{source}_4)\wedge d\rho\wedge d\phi$.
The circle-first quotient therefore has
$\operatorname{or}(Y)=-\operatorname{or}(\text{source}_4)
\wedge d\log\rho$. In particular its identity with the
weighted positive-circle product reverses that product
orientation. This observation alone does not assign the sign
of the return, since it also changes the normal meridian.

Let $\kappa$ be the actual sign in
$\partial V=\kappa E\times S^1_\theta$. The compression
has boundary $F-E$, so the glued trace has
$\partial X=\kappa F\times S^1_\theta$. In its outer collar,
\[
 \Omega_Z=\Omega_F\wedge\Omega_{\nu,\mathrm{Hopf}}
                  \wedge dr\wedge d\theta,
 \qquad
 \Omega_X=\kappa\Omega_F\wedge dr\wedge d\theta.
\]
The oriented normal sphere is consequently $\kappa$ times
the Hopf meridian, and the extended actual Euler class has
$\pi_*x=\kappa$. Put $\widehat x=\kappa x$ and
$e_v=2\widehat x-\pi^*b$. The framed graph section gives
$b|_V=2\kappa c_1(M)$, with the same retained boundary
correction as P3K.5. Therefore $b^2=4M_{\rm rel}^2$.
Fibre integration and both outer/Jin orientation reversals
of P3K.7 give
\[
 A(\Phi)=\kappa\frac{3\langle b^2,[\bar X]\rangle
                   -\langle p_1(T\bar X),[\bar X]\rangle}{24}.
                                                               \tag{K24.10}
\]
In the actual affine page the complex boundary of a large
$a$-disc is $+\arg a$. Formula (K24.6) identifies it with
outward gamma, and (K24.5) identifies the positive parameter
with that same outward gamma. Hence $\kappa=+1$. The actual
line, Thom sign and parameter loop are now all fixed. Equations
(K24.9)--(K24.10) give
\[
 A(\Phi)=\frac{12(-2m)+24\chi}{24}
        =\chi-m=-(3l+2).
\]
The Hopf transfer P3F.22 gives (K24.1). The residues
$26-3l$ exhaust $\mathbb Z/28$ because $\gcd(3,28)=1$;
the mapping-torus assertion follows from P3F.24.

\par\medskip\noindent\textbf{The explicit III*$+$II$+I_1$ source.}
For $y^2=x^3+tx-t$ the discriminant is a nonzero constant
times $t^2(4t+27)$. At zero its coefficient orders are
$(1,1,2)$, giving II, and $-27/4$ gives $I_1$. At infinity,
\[
 y_0^2=x_0^3+q^3x_0-q^5,\qquad G=(q^2,q^3),
\]
giving III*. The source is rational by
$t=(y^2-x^3)/(x-1)$. Its fourth good root sends $G$ to the
nonzero fixed two-torsion point, so its $E_7$ component,
height $1/2$, primitivity and absence of torsion are the same
marked data as in P3F.37. In particular $2G\cdot O=0$ and
its Hom line for $Q=G$ is the whole-fibre line of degree one.
Use its sole $I_1$ as the zero. The actual reduced root is
\[
 Y^2=X^3+zX-z^{2l+3},\qquad
 G'=(z^{2l+2},z^{3l+3}).
\]
The quadratic part is again $Y^2-zX$; the same resolution,
deck, unique $m$-division, line norm and meridian construct its
III cap. The basis $O,F,E_7,G$ has the identical unimodular
Gram matrix and $TM=M-F$. The two vanishing cycles of its
actual II fibre meet once and generate the elliptic lattice.
This gives the source-complement simple connectivity and every
unchanged integral gluing map in P3F.37. Thus this source too
gives the stated proper family of standard six-spheres.

Its actual page has one III, $m$ II and $m$ nodal fibres,
again of Euler number $3+3m$. The same lifted vanishing
relations prove (K24.8) relative to the binding. Its doubling
slope is $(1+3z^{4l+3})/(2z^{3l+2})$, with the same leading
zero-section parameter. Every line, signature and oriented
trace calculation above therefore has the same value. The
independent certificates are
\path{code/part3/reduced_iii_return.py} and
\path{code/part3_root/reduced_iii_return_audit.py}; the working
note retains the superseded character convention verbatim in
its dated historical record.
\end{proof}
